\documentclass[11pt]{article}
\usepackage[a4paper,margin=1in]{geometry}
\usepackage{amsmath,amssymb,amsthm}
\usepackage{graphicx}
\usepackage{hyperref}

\newtheorem{theorem}{Theorem}

\newcommand{\K}{\mathbb K}
\newcommand{\C}{\mathbb C}
\newcommand{\R}{\mathbb R}
\newcommand{\Pol}{\mathcal P}
\newcommand{\Lin}{\mathcal L}
\newcommand{\norm}[1]{\left\lVert #1 \right\rVert}

\title{A scalar counterexample for nonlinear polynomial Daugavet centers}
\author{Run \texttt{fa\_banach\_001}, agent lane 12}
\date{2026-06-20}

\begin{document}
\maketitle

\section*{Source and question}

The source paper is Sheldon Dantas, Miguel Martin, and Yoel Perreau,
\emph{The Daugavet property is equivalent to the polynomial Daugavet property},
arXiv:2507.06143 \cite{DMP}.

After proving that every linear Daugavet center is a polynomial Daugavet center,
the authors ask on page 11 whether the linearity assumption can be removed:
\begin{quote}\small
Given a non-null polynomial \(Q\in\Pol(X,Y)\) satisfying
\[
\norm{Q+T}=\norm{Q}+\norm{T}
\]
for every rank-one \(T\in\Lin(X,Y)\), must \(Q\) be a polynomial Daugavet
center, i.e. must the same equality hold when \(T\) is a rank-one polynomial
from \(X\) to \(Y\)?
\end{quote}

\begin{center}
\includegraphics[width=0.93\linewidth]{figures/open_problem_crop.png}
\end{center}

The answer is negative, even in dimension one.

\section*{Definitions Used}

The source paper defines a polynomial \(P:X\to Y\) as a finite linear
combination of homogeneous polynomials, including constant maps as
\(0\)-homogeneous polynomials.  It defines a rank-one polynomial as a map
\[
P(x)=p(x)y_0
\]
with \(p\in\Pol(X)\) scalar-valued and \(y_0\in Y\).  A non-null polynomial
\(Q\in\Pol(X,Y)\) is a polynomial Daugavet center if
\[
\norm{Q+P}=\norm{Q}+\norm{P}
\]
for every rank-one polynomial \(P\in\Pol(X,Y)\).

\section*{Proof Intuition}

Linear perturbations of the scalar polynomial \(Q(z)=z^2\) are too symmetric to
detect the obstruction.  If \(T(z)=az\), the maximum of \(|z^2+az|\) on the
unit ball is attained at a boundary point whose phase aligns \(z\) with \(a\),
giving exactly \(1+|a|\).  Polynomial perturbations are much more flexible:
the rank-one polynomial \(P=-Q\) cancels \(Q\) identically.  Thus the linear
Daugavet-center test is strictly weaker than the polynomial one.

\begin{theorem}
Over either \(\K=\R\) or \(\K=\C\), there are one-dimensional Banach spaces
\(X=Y=\K\) and a non-null polynomial \(Q\in\Pol(X,Y)\) such that
\[
\norm{Q+T}=\norm{Q}+\norm{T}
\]
for every rank-one linear operator \(T\in\Lin(X,Y)\), but \(Q\) is not a
polynomial Daugavet center.
\end{theorem}

\begin{proof}
Take \(X=Y=\K\) with the usual absolute-value norm and define
\[
Q(z)=z^2 \qquad (z\in\K).
\]
Then \(\norm{Q}=1\).  Every linear operator \(T:X\to Y\) has the form
\(T(z)=az\) for some \(a\in\K\), and in dimension one every nonzero such
operator has rank one.  The zero operator also causes no issue.

For \(|z|\leq 1\), the triangle inequality gives
\[
|Q(z)+T(z)|=|z^2+az|\leq |z|^2+|a||z|\leq 1+|a|.
\]
Conversely, if \(a\neq 0\), choose \(z=a/|a|\) in the complex case and
\(z=\operatorname{sign}(a)\) in the real case.  Then \(|z|=1\) and
\[
|z^2+az|=|z|\,|z+a|=1+|a|.
\]
If \(a=0\), the same equality is clear.  Hence
\[
\norm{Q+T}=1+|a|=\norm{Q}+\norm{T}
\]
for every rank-one linear \(T\).

However, \(P=-Q\) is a rank-one polynomial from \(X\) to \(Y\): indeed
\[
P(z)=p(z)y_0,\qquad p(z)=z^2,\quad y_0=-1.
\]
It has \(\norm{P}=1\), while
\[
\norm{Q+P}=0 < 2=\norm{Q}+\norm{P}.
\]
Therefore \(Q\) is not a polynomial Daugavet center.
\end{proof}

\section*{Verification Notes}

The proof is exact.  The accompanying script
\texttt{code/verify\_scalar\_example.py} samples the complex unit circle for
several values of \(a\) and illustrates
\(\sup_{|z|\leq1}|z^2+az|=1+|a|\), and records the failure for \(P=-Q\).

\section*{Novelty Check and Limitations}

On 2026-06-20, I searched for exact phrases including
\begin{quote}\small
\texttt{"polynomial Daugavet center" "rank-one" "linear"},\\
\texttt{"we do not know" "polynomial Daugavet center"},
\end{quote}
and the source title.  The only direct hit found was the source arXiv paper.
This counterexample answers only the nonlinear polynomial-center question; it
does not address the separate Bourgain-type lemma question for the weak
polynomial topology or the open problem of whether DPr implies WODP.

\section*{Human Review Recommendation}

This should be reviewed as a likely valid counterexample.  The key point is
whether the source definition of rank-one polynomial indeed includes
\(P(x)=p(x)y_0\) for arbitrary scalar polynomial \(p\in\Pol(X)\).  Under that
definition, \(P=-Q\) is rank-one and the counterexample is immediate.

\begin{thebibliography}{9}
\bibitem{DMP}
Sheldon Dantas, Miguel Martin, and Yoel Perreau,
\emph{The Daugavet property is equivalent to the polynomial Daugavet property},
arXiv:2507.06143, 2025.
\end{thebibliography}

\end{document}
